\chapter{Recursos utilizados}
\label{recursosUtilizados}

Este capítulo resume los recursos humanos, software y hardware empleados para el análisis, diseño, implementación, validación y documentación de \textit{Nuvlo}. La finalidad es dejar constancia del entorno real utilizado y facilitar que el proyecto pueda reproducirse o mantenerse posteriormente.

\section{Recursos humanos}

Los recursos humanos implicados en el proyecto se organizan en torno al desarrollo, la dirección académica y la supervisión del trabajo.

\begin{table}[htbp]
\centering
\small
\begin{tabularx}{\textwidth}{p{0.28\textwidth} X}
\toprule
\textbf{Recurso} & \textbf{Participación} \\
\midrule
\textit{Autor del TFG} & Responsable del análisis, diseño, implementación, pruebas, documentación y preparación de la memoria del \textit{TFG}. \\
\midrule
\textit{Tutores del TFG} & Orientación académica, revisión del enfoque del \textit{TFG}, seguimiento de la memoria y supervisión general del trabajo. \\
\bottomrule
\end{tabularx}
\caption{Recursos humanos utilizados.}
\label{tab:recursosHumanos}
\end{table}

\section{Recursos software}

La Tabla~\ref{tab:recursosSoftware} recoge las herramientas software principales utilizadas durante el desarrollo. Se incluyen tanto tecnologías de implementación como herramientas de documentación, control de versiones y validación.

\begin{table}[htbp]
\centering
\small
\begin{tabularx}{\textwidth}{p{0.28\textwidth} p{0.24\textwidth} X}
\toprule
\textbf{Recurso} & \textbf{Tipo} & \textbf{Uso principal} \\
\midrule
\textit{JavaScript} / \textit{Node.js} & Lenguaje y entorno de ejecución & Implementación del servidor, scripts de desarrollo y lógica de aplicación~\cite{NodeDocs}. \\
\midrule
\textit{React} & Biblioteca de interfaz & Construcción del cliente web mediante componentes reutilizables~\cite{ReactDocs}. \\
\midrule
\textit{Vite} & Herramienta de desarrollo & Servidor de desarrollo y empaquetado del \textit{frontend}~\cite{ViteDocs}. \\
\midrule
\textit{Express} & Framework web & Definición de rutas HTTP, \textit{middlewares} y \textit{API} del \textit{backend}~\cite{ExpressDocs}. \\
\midrule
\textit{Prisma} & \textit{ORM} & Definición del modelo de datos, migraciones y acceso a \textit{PostgreSQL}~\cite{PrismaDocs}. \\
\midrule
\textit{PostgreSQL} & Base de datos relacional & Persistencia de usuarios, proyectos, incidencias, métricas, alertas, actividad y datos sincronizados~\cite{PostgreSQLDocs}. \\
\midrule
\textit{Redis} & Almacén temporal & Gestión de cachés de corta duración y estado de sincronización con \textit{TTL}~\cite{RedisDocs}. \\
\midrule
\textit{Docker Compose} & Orquestación local & Ejecución reproducible de servicios auxiliares como \textit{PostgreSQL} y \textit{Redis}~\cite{DockerComposeDocs}. \\
\midrule
Git y \textit{GitHub} & Control de versiones & Historial de cambios, repositorio remoto y flujo de integración continua. \\
\midrule
\textit{GitHub Actions} & Integración continua & Compilación automatizada de la memoria \textit{LaTeX} y validación básica del repositorio. \\
\midrule
\textit{Vitest} & Pruebas automatizadas & Pruebas unitarias e integración de lógica de métricas, \textit{API} y servicios~\cite{VitestGuide}. \\
\midrule
\textit{Playwright} & Pruebas \textit{E2E} & Validación de flujos principales de la interfaz en navegador~\cite{PlaywrightDocs}. \\
\midrule
\textit{LaTeX} & Documentación académica & Elaboración y compilación de la memoria del proyecto. \\
\midrule
Draw.io / diagrams.net & Diagramación & Creación de diagramas de casos de uso, secuencia, arquitectura, base de datos y vistas. \\
\midrule
Visual Studio Code & Entorno de desarrollo & Edición del código fuente, documentación y ejecución de terminales integradas. \\
\bottomrule
\end{tabularx}
\caption{Recursos software utilizados.}
\label{tab:recursosSoftware}
\end{table}

\section{Recursos hardware}

El desarrollo se ha realizado en un equipo personal con sistema \textit{Windows} y entorno \textit{Linux} mediante \textit{WSL}, siglas de \textit{Windows Subsystem for Linux}. Esta configuración permite trabajar con herramientas habituales de desarrollo web y, al mismo tiempo, mantener compatibilidad con comandos y entornos \textit{Linux} utilizados en servidores y sistemas de integración continua.

\begin{table}[htbp]
\centering
\small
\begin{tabularx}{\textwidth}{p{0.30\textwidth} X}
\toprule
\textbf{Recurso} & \textbf{Uso principal} \\
\midrule
Ordenador portátil personal & Desarrollo de la aplicación, ejecución local del \textit{frontend}, \textit{backend}, base de datos, \textit{Redis}, pruebas y compilación de la memoria. \\
\midrule
\textit{Windows} 11 con \textit{WSL} & Entorno de trabajo mixto para utilizar herramientas de escritorio y entorno \textit{Linux} de desarrollo. \\
\midrule
\textit{Docker} Desktop & Ejecución local de contenedores y conexión con \textit{WSL}. \\
\midrule
Navegador web & Pruebas funcionales, revisión visual de la interfaz y validación del flujo \textit{OAuth} con \textit{Atlassian}. \\
\midrule
\textit{GitHub Actions} & Ejecución remota de tareas de integración continua sobre infraestructura de \textit{GitHub}. \\
\bottomrule
\end{tabularx}
\caption{Recursos hardware y entorno de ejecución utilizados.}
\label{tab:recursosHardware}
\end{table}

\section{Servicios externos}

Además del entorno local, el proyecto utiliza servicios externos necesarios para la integración y validación de la aplicación.

\begin{itemize}
    \item \textbf{\textit{Jira Cloud}}: fuente principal de datos para proyectos, incidencias, tableros, \textit{Sprints} y transiciones.
    \item \textbf{\textit{Atlassian Developer Console}}: configuración de la aplicación \textit{OAuth} 2.0 3LO y permisos de acceso.
    \item \textbf{\textit{GitHub}}: alojamiento del repositorio, control remoto del código y ejecución de flujos de integración continua.
\end{itemize}

Estos servicios condicionan parte de la arquitectura del sistema, especialmente la autenticación, la sincronización de datos y la necesidad de mantener una demo \textit{offline} que permita comprobar la aplicación aunque una \textit{API} externa no esté disponible.
